% ==========================================
% 1. Introduction
% ==========================================
%
% ------------------------------------------
% Status of this section
% ------------------------------------------
% Semi-stable: parts 1-3 are written now; parts 4-5 are written last.
% Length: about 1.5-2 pages. Broad and short; the details belong to
% the problem definition (section 2).

\section{Introduction}
\label{sec:introduction}

% ------------------------------------------
% 1. Context and motivation
% ------------------------------------------
% TODO [NOW]: We present the domain and what Splinter3D is, for a reader who
%             discovers the topic. We explain why the topic matters
%             (applications, needs, stakes), with references for each claim.

% ------------------------------------------
% 2. Problem statement
% ------------------------------------------
% TODO [NOW]: We state the problem in a few sentences and point to
%             Section~\ref{sec:problem-definition} for the formal definition.
%             We do not duplicate the problem definition here.
% TODO [NOW]: We state, in one or two sentences, what existing work does not solve
%             (details are given in Section~\ref{sec:state-of-the-art}).

% ------------------------------------------
% 3. Objectives and research questions
% ------------------------------------------
% TODO [NOW]: We state the goal of the study and list the research questions
%             (RQ1, RQ2, ...), worded exactly as in the problem definition.
% TODO [NOW]: We state the scope in one sentence: what the study covers and
%             what it deliberately leaves out.

% ------------------------------------------
% 4. Contributions
% ------------------------------------------
% TODO [LAST]: We list the contributions as a short itemized list. Each
%              contribution is verifiable in the report (it points to a section,
%              an experiment or a result) and none is announced without
%              a corresponding result.
% Example structure:
%   \begin{itemize}
%       \item We ...  (see Section~\ref{sec:results}).
%       \item We ...
%   \end{itemize}

% ------------------------------------------
% 5. Report outline
% ------------------------------------------
% TODO [LAST]: We describe the organization of the report in one short
%              paragraph, using \ref for every section (never hard-coded numbers).
% TODO [CHECK]: The outline matches the sections that actually exist.
% TODO [CHECK]: The objectives and contributions are still consistent with the
%               final results; if the project changed direction, we rewrite them.
